%% toc: P15 171 171.5,171.6,171.7,171.8
\section{Uniform convergence without a countable core, bounded losses and lower bounds: complete proofs}
\label{sec:annex-p15-w02b-p15-065-beyond-core}

Notation follows
\cref{def:p15-vcl-event-class,def:p15-w02b-p15-065-outer,def:p15-w02b-p15-065-sym-measurable,def:p15-w02b-p15-065-level-class,def:p15-w02b-p15-065-learner,def:p15-w02b-p15-065-noise-conditions}.
For \(\sigma\in\{-1,1\}^m\), \(T_\sigma:Z^{2m}\to Z^{2m}\) exchanges the
coordinates \(z_i\) and \(z'_i\) for each \(i\) with \(\sigma_i=-1\) and
leaves the others in place.

\begin{annexitem}{annex:p15-w02b-p15-065-outer-fubini}{thm:p15-w02b-p15-065-outer-fubini}{Measurable majorants and their sections}
\begin{proof}
\emph{Elementary properties.} If \(f\leq f'\), every measurable majorant of
\(f'\) majorizes \(f\), so \(E^{*}f\leq E^{*}f'\). For \(c>0\), \(g\geq f\)
iff \(cg\geq cf\), so \(E^{*}(cf)=cE^{*}f\). If \(g\geq\mathbf1_A\) is
measurable, then \(\{g\geq1\}\) is a measurable superset of \(A\) and
\(\int g\geq Q\{g\geq1\}\geq Q^{*}(A)\); the indicator of any measurable
superset is a majorant, so \(E^{*}\mathbf1_A=Q^{*}(A)\). If \(h\leq f\) is
measurable, every majorant \(g\) of \(f\) satisfies \(\int g\geq\int h\).
For sets \(A_n\) and \(\eta>0\), choose measurable \(B_n\supseteq A_n\) with
\(Q(B_n)\leq Q^{*}(A_n)+\eta2^{-n}\); then \(\bigcup_nB_n\supseteq\bigcup_nA_n\)
gives \(Q^{*}(\bigcup_nA_n)\leq\sum_nQ^{*}(A_n)+\eta\), and \(\eta\) is
arbitrary. If \(B_1\subseteq A\subseteq B_2\) with
\(Q(B_2\setminus B_1)=0\), then \(Q^{*}(A)\leq Q(B_2)\), and every measurable
\(B\supseteq A\) contains \(B_1\), so \(Q(B)\geq Q(B_1)=Q(B_2)\). Hence
\(Q^{*}(A)=Q(B_2)=\overline Q(A)\).

\emph{The product inequality.} Fix \(\eta>0\) and a
\(\mathcal F_1\otimes\mathcal F_2\)-measurable \(g\geq f\) with
\(\int g\,d(Q_1\otimes Q_2)\leq E^{*}f+\eta\). Replacing \(g\) by
\(\min\{g,1\}\) keeps it measurable, keeps \(g\geq f\) because \(f\leq1\),
and does not increase the integral; so \(0\leq g\leq1\). For each
\(\omega_1\), the section \(g(\omega_1,\cdot)\) is
\(\mathcal F_2\)-measurable, since sections of product-measurable functions
are measurable, and it majorizes \(f(\omega_1,\cdot)\). Therefore
\[
 E_{Q_2}^{*}f(\omega_1,\cdot)\leq G(\omega_1):=\int g(\omega_1,\omega_2)\,Q_2(d\omega_2).
\]
By Tonelli's theorem \(G\) is \(\mathcal F_1\)-measurable and
\(\int G\,dQ_1=\int g\,d(Q_1\otimes Q_2)\). So \(G\) is a measurable majorant
of \(\omega_1\mapsto E_{Q_2}^{*}f(\omega_1,\cdot)\), and
\[
 E_{Q_1}^{*}\bigl[E_{Q_2}^{*}f(\omega_1,\cdot)\bigr]\leq\int G\,dQ_1
 \leq E^{*}f+\eta .
\]
Let \(\eta\to0\). The failure of the reverse inequality under the continuum
hypothesis is item~4 of \cref{thm:p15-w02b-p15-065-ch-failure}, proved in
\cref{annex:p15-w02b-p15-065-ch-failure}. There each section of the set has
outer probability zero, so the iterated outer expectation of its indicator
is zero, while its joint outer probability is one.
\end{proof}
\end{annexitem}

\begin{annexitem}{annex:p15-w02b-p15-065-core-sufficient}{thm:p15-w02b-p15-065-core-sufficient}{Core approximation on a pooled sample}
\begin{proof}
Fix \(w\in Z^{2m}\) and \(C\in\mathcal C\), with core events \(C_j\) whose
indicators converge pointwise to \(\mathbf1_C\). Binary sequences converge
only by becoming eventually constant, so at each of the \(2m\) coordinates
of \(w\) the indicators eventually agree with \(\mathbf1_C\). Hence for all
large \(j\), \(P_m(C_j)=P_m(C)\) and \(P'_m(C_j)=P'_m(C)\). Consequently
\(|P_m(C)-P'_m(C)|\leq\sup_{C'\in\mathcal C_0}|P_m(C')-P'_m(C')|\), and the
reverse inequality holds because \(\mathcal C_0\subseteq\mathcal C\). So
\(D_m\) equals the supremum over the core. If \(P_m(C)=0\), the same late
core events satisfy \(P_m(C_j)=0\) and \(P'_m(C_j)=P'_m(C)\). Therefore
\[
 K_m(w)=\sup\Bigl(\{0\}\cup\{P'_m(C_j)\mathbf1\{P_m(C_j)=0\}:C_j\in\mathcal C_0\}\Bigr),
\]
where a term with \(P_m(C_j)\ne0\) contributes \(0\), already in the set.
For \(C\in\mathcal A\), \(w\mapsto P_m(C)\) and \(w\mapsto P'_m(C)\) are
finite sums of measurable indicators. A countable supremum of measurable
functions is measurable. Both \(D_m\) and \(K_m\) are thus
\(\mathcal A^{\otimes2m}\)-measurable, hence measurable for the completion
under any product law.
\end{proof}
\end{annexitem}

\begin{annexitem}{annex:p15-w02b-p15-065-outer-uniform}{thm:p15-w02b-p15-065-outer-uniform}{Sample-by-sample witnesses, invariant swaps and a finite union}
\begin{proof}
\emph{Swaps preserve the completed law.} Each \(T_\sigma\) is a
permutation of coordinates, is its own inverse, and is measurable for
\(\mathcal A^{\otimes2m}\). The image of \(P^{\otimes2m}\) under
\(T_\sigma\) agrees with \(P^{\otimes2m}\) on measurable rectangles
\(A_1\times\cdots\times A_{2m}\): exchanging two coordinates exchanges two
factors \(P(A_i)\) and \(P(A_{m+i})\) of the same product. Rectangles form a
\(\pi\)-system generating the product \(\sigma\)-algebra, so the two laws
coincide. If \(N\subseteq M\) with \(M\) measurable and null, then
\(T_\sigma N\subseteq T_\sigma M\), again measurable and null. Hence
\(T_\sigma\) maps completion-measurable sets to completion-measurable sets
and \(\overline{P^{\otimes2m}}(T_\sigma^{-1}A)=\overline{P^{\otimes2m}}(A)\).

\emph{Fixed-pool bounds.} Fix \(w\). For \(C\in\mathcal C\) put
\(a_i(C)=\mathbf1_C(z_i)-\mathbf1_C(z'_i)\in\{-1,0,1\}\). The first half of
\(T_\sigma w\) holds \(z_i\) where \(\sigma_i=1\) and \(z'_i\) where
\(\sigma_i=-1\), so
\(P_m(C)-P'_m(C)\) evaluated at \(T_\sigma w\) equals
\(m^{-1}\sum_i\sigma_ia_i(C)\). The vector \(a(C)\), and the pair of counts
of \(C\) in the two halves of \(T_\sigma w\) for each \(\sigma\), depend
only on the trace of \(C\) on the ordered \(2m\)-tuple \(w\). At most
\(\Pi_{\mathcal C}(2m)\) traces occur. Let \(\sigma\) be uniform on
\(\{-1,1\}^m\).

For one vector \(a\in\{-1,0,1\}^m\) with \(s=\sum_ia_i^2\), the event
\(|\sum_i\sigma_ia_i|>m\varepsilon/2\) is empty if \(s=0\). Otherwise, for
\(\lambda>0\), the series \(\cosh u=\sum_ku^{2k}/(2k)!\) and
\((2k)!\geq2^kk!\) give \(\cosh u\leq e^{u^2/2}\). Independence of the signs
and Markov's inequality give, for \(t=m\varepsilon/2\),
\[
 \Pr_\sigma\Bigl(\sum_i\sigma_ia_i>t\Bigr)\leq e^{-\lambda t}\prod_i\cosh(\lambda a_i)
 \leq e^{-\lambda t+\lambda^2s/2}.
\]
At \(\lambda=t/s\) this is \(e^{-t^2/(2s)}\leq e^{-t^2/(2m)}=e^{-m\varepsilon^2/8}\)
because \(s\leq m\). The negative tail is the same by the symmetry
\(\sigma\mapsto-\sigma\). A union over the traces gives
\[
 \Pr_\sigma\bigl(D_m(T_\sigma w)>\varepsilon/2\bigr)\leq2\Pi_{\mathcal C}(2m)e^{-m\varepsilon^2/8}.
 \tag{F1}
\]
For the second bound, fix a trace. If some pair has both points in \(C\),
the first half of \(T_\sigma w\) always meets \(C\). Otherwise let \(k\)
pairs have exactly one point in \(C\). The first half avoids \(C\) iff each
of these \(k\) points is sent to the second half, which has probability
\(2^{-k}\), and then the second half contains exactly \(k\) of them. The
requirement \(k/m>\varepsilon/2\) forces \(2^{-k}<2^{-m\varepsilon/2}\).
A union over the traces gives
\[
 \Pr_\sigma\bigl(K_m(T_\sigma w)>\varepsilon/2\bigr)\leq\Pi_{\mathcal C}(2m)2^{-m\varepsilon/2}.
 \tag{F2}
\]
Both bounds hold for every \(w\), with no measurability needed, because
\(\sigma\) ranges over a finite set.

\emph{Item~1.} If \(m\varepsilon^2<2\), then
\(4\Pi_{\mathcal C}(2m)e^{-m\varepsilon^2/8}\geq4e^{-1/4}>1\) and the bound
one is trivial. Let \(m\varepsilon^2\geq2\). For each \(z\in Z^m\) with
\(\Delta_m(z)>\varepsilon\), choose \(C_z\in\mathcal C\) with
\(|P_m(C_z)-P(C_z)|>\varepsilon\); the choice is made separately for each
\(z\), and no measurability of \(z\mapsto C_z\) is used. For this fixed
\(C_z\), the set
\(G_z=\{z':|P'_m(C_z)-P(C_z)|\leq\varepsilon/2\}\) is
\(\mathcal A^{\otimes m}\)-measurable. Since \(mP'_m(C_z)\) is binomial,
\(P'_m(C_z)\) has mean \(P(C_z)\) and variance at most \(1/(4m)\), and Chebyshev's
inequality gives \(P^{\otimes m}(G_z)\geq1-1/(m\varepsilon^2)\geq1/2\). On
\(G_z\),
\(|P_m(C_z)-P'_m(C_z)|\geq|P_m(C_z)-P(C_z)|-|P'_m(C_z)-P(C_z)|>\varepsilon/2\),
so \(f=\mathbf1\{D_m>\varepsilon/2\}\) satisfies \(f(z,\cdot)\geq\mathbf1_{G_z}\).
By the elementary properties in
\cref{annex:p15-w02b-p15-065-outer-fubini},
\(E^{*}f(z,\cdot)\geq1/2\). Thus
\(\mathbf1\{\Delta_m(z)>\varepsilon\}\leq2E^{*}f(z,\cdot)\) for every \(z\),
and monotonicity, homogeneity and the product inequality give
\[
 P^{*}(\Delta_m>\varepsilon)\leq2E^{*}_{P^{\otimes m}}\bigl[E^{*}_{P^{\otimes m}}f(z,\cdot)\bigr]
 \leq2E^{*}_{P^{\otimes2m}}f=2\,\overline{P^{\otimes2m}}(A),
\]
where \(A=\{D_m>\varepsilon/2\}\) is completion-measurable by hypothesis.
By swap invariance and linearity of the integral,
\[
 \overline{P^{\otimes2m}}(A)=2^{-m}\sum_\sigma\overline{P^{\otimes2m}}(T_\sigma^{-1}A)
 =\int\Pr_\sigma\bigl(D_m(T_\sigma w)>\varepsilon/2\bigr)\,\overline{P^{\otimes2m}}(dw),
\]
a finite sum of integrals of completion-measurable indicators. By (F1) the
integrand is at most \(2\Pi_{\mathcal C}(2m)e^{-m\varepsilon^2/8}\), which
proves item~1. This item uses only the measurability of \(D_m\) for the
given \(m\).

\emph{Item~2.} Let \(m\varepsilon\geq8\). For \(z\) in the bad event, choose
\(C_z\) with \(P_m(C_z)=0\) and \(p=P(C_z)>\varepsilon\). The ghost count
\(B=mP'_m(C_z)\) is binomial with mean \(mp\) and variance at most \(mp\);
Chebyshev's inequality gives
\(P(B<mp/2)\leq4/(mp)<4/(m\varepsilon)\leq1/2\). On \(\{B\geq mp/2\}\),
which is measurable in \(z'\),
\(K_m(z,z')\geq P'_m(C_z)\geq p/2>\varepsilon/2\). The same chain with
\(f=\mathbf1\{K_m>\varepsilon/2\}\), its completion-measurability, swap
invariance and (F2) gives the bound \(2\Pi_{\mathcal C}(2m)2^{-m\varepsilon/2}\).

\emph{Reduction.} Under a countable core, \(\Delta_m\) and the bad
consistent event are measurable by \cref{thm:p15-vcl-core}, and outer
probability of a measurable set is its probability.
\end{proof}
\end{annexitem}

\begin{annexitem}{annex:p15-w02b-p15-065-outer-learning}{thm:p15-w02b-p15-065-outer-learning}{Inclusions of bad events and the arithmetic of sufficient sizes}
\begin{proof}
\emph{Growth.} On a labelled tuple \(((x_i,y_i))_{i\leq t}\), the loss
vector of \(h\) is its prediction vector on \((x_i)\) combined with the
fixed label vector by exclusive or. This is a bijection of \(\{0,1\}^t\),
so the number of loss traces equals the number of prediction traces, at
most \(\Pi_H(t)\). With every label zero the two families coincide, so the
maximum over labelled tuples is exactly \(\Pi_H(t)\).

\emph{Item~1.} Each empirical risk lies in \(\{0,1/m,\ldots,1\}\); the
attained values form a nonempty finite set, whose minimum is attained by
some member. Hence exact and \(\eta\)-approximate minimizers exist on every
sample; the map \(\widehat h\) is any choice of one. Write
\(\Delta_m(\mathcal E_H)=\sup_h|\widehat R_m(h)-R(h)|\). On
\(\{\Delta_m(\mathcal E_H)\leq\varepsilon\}\), for every \(h\in H\),
\(R(\widehat h)\leq\widehat R_m(\widehat h)+\varepsilon
\leq\widehat R_m(h)+\eta+\varepsilon\leq R(h)+2\varepsilon+\eta\), and the
infimum over \(h\) gives \(R(\widehat h)\leq R_H^{*}+2\varepsilon+\eta\).
So the bad set is contained in \(\{\Delta_m(\mathcal E_H)>\varepsilon\}\),
and monotonicity of \(P^{*}\) with item~1 of
\cref{thm:p15-w02b-p15-065-outer-uniform} applied to \(\mathcal E_H\)
proves the bound.

\emph{Item~2.} The set \(N=\bigcup_{i\leq m}\{h^{*}(X_i)\ne Y_i\}\) is
measurable with probability at most \(mR(h^{*})=0\). Off \(N\), \(h^{*}\)
has no training error, so \(\widehat h\) has none. If also
\(R(\widehat h)>\varepsilon\), then the loss event of \(\widehat h\) has
empirical frequency zero and probability above \(\varepsilon\). So the bad
set lies in \(N\) together with the event of item~2 of
\cref{thm:p15-w02b-p15-065-outer-uniform} for \(\mathcal E_H\).
Subadditivity of \(P^{*}\) gives the bound.

\emph{Item~3.} \Cref{annex:p15.vcl.samples} shows, using only Sauer's bound
\cref{thm:p15-vcl-sauer}, that for every \(m\) at least the first displayed
size \(4\Pi(2m)e^{-m\varepsilon^2/8}\leq\delta\), and for every \(m\) at
least the second \(2\Pi(2m)2^{-m\varepsilon/2}\leq\delta\). These are the
right-hand sides of items~1 and~2, so the same sizes suffice.
\end{proof}
\end{annexitem}

\begin{annexitem}{annex:p15-w02b-p15-065-singletons}{thm:p15-w02b-p15-065-singletons}{Counting coincidences}
\begin{proof}
\emph{Item~1.} On an ordered \(t\)-tuple the traces are the empty set and
the sets of positions holding one fixed value. There are at most \(t+1\),
and exactly \(t+1\) when the points are distinct, which is possible since
\([0,1]\) is infinite. A single point is shattered by \(\{x\}\) and
\(\varnothing\). Two distinct points are not, since no member contains
both. If \(\mathcal C_0\) is a countable subfamily, the union of its members
is countable. Pick \(x\) outside it. Every indicator from \(\mathcal C_0\)
is zero at \(x\), so none converges to \(\mathbf1_{\{x\}}(x)=1\).

\emph{Item~2.} For \(w\in[0,1]^{2m}\) let
\(N_1(x)=\#\{i\leq m:z_i=x\}\) and \(N_2(x)=\#\{i\leq m:z'_i=x\}\). For
\(C=\{x\}\) the two-sample difference is \((N_1(x)-N_2(x))/m\), which is
zero unless \(x\) is a coordinate of \(w\); \(\varnothing\) contributes
zero. Writing \(w_k\) for the coordinates,
\begin{align*}
 D_m(w)&=\max_{k\leq2m}\frac{|N_1(w_k)-N_2(w_k)|}m,\\
 K_m(w)&=\max\Bigl\{0,\max_{k\leq2m}\frac{N_2(w_k)}m\mathbf1\{N_1(w_k)=0\}\Bigr\}.
\end{align*}
Each \(N_r(w_k)\) is a sum of indicators of the closed sets
\(\{w_i=w_k\}\), so both functions are Borel. For the one-sample quantity,
let \((a_j)\) list the countably many atoms of \(P\). Then
\begin{align*}
 \Delta_m(z)=\max\Bigl\{&\max_{i\leq m}\Bigl|\frac{N_1(z_i)}m-\sum_jP\{a_j\}\mathbf1\{z_i=a_j\}\Bigr|,\\
 &\sup_jP\{a_j\}\mathbf1\{a_j\notin\{z_1,\ldots,z_m\}\}\Bigr\},
\end{align*}
because a non-atom outside the sample contributes zero; this is Borel. The
growth at \(2m\) is \(2m+1\) by item~1.

\emph{Item~3.} Under the uniform law each diagonal \(\{z_i=z_k\}\),
\(i\ne k\), is null, so almost surely the sample points are distinct and
there are no atoms. Then \(\Delta_m=\max_iN_1(z_i)/m=1/m\). Under the point
mass at \(a\), almost surely every \(z_i=a\). Then \(P_m\{a\}=1=P\{a\}\),
and every other member has empirical and true probability zero.
\end{proof}
\end{annexitem}

\begin{annexitem}{annex:p15-w02b-p15-065-ch-failure}{thm:p15-w02b-p15-065-ch-failure}{A well-ordering with countable initial segments}
\begin{proof}
Under the continuum hypothesis \([0,1]\) has cardinality \(\aleph_1\), so
there is a bijection \(\varphi\) from \([0,1]\) onto the set \(\omega_1\) of
countable ordinals. Put \(x\prec y\) iff \(\varphi(x)<\varphi(y)\). This is
a well-ordering, and \(I_x=\varphi^{-1}(\{\alpha:\alpha\leq\varphi(x)\})\) is
countable because \(\varphi(x)\) is a countable ordinal.

\emph{Item~1.} A countable set is a countable union of closed singletons,
hence Borel, and has Lebesgue measure zero. The class is a chain:
\(x\preceq x'\) implies \(I_x\subseteq I_{x'}\). If \(a\prec b\), every
\(I_x\) containing \(b\) contains \(a\), so no two-point set is shattered. A
point \(a\) that is not \(\prec\)-least is shattered: \(I_a\) contains it
and \(I_x\) with \(x\prec a\) does not. The traces of a chain on a
\(t\)-tuple form a chain of subsets of \(t\) positions, so there are at
most \(t+1\) of them.

\emph{Item~2.} Given \(z_1,\ldots,z_m\), let \(x\) be the \(\prec\)-largest
of these finitely many points. Then \(I_x\) contains every \(z_i\), so
\(P_m(I_x)=1\), while \(P(I_x)=0\). Every member has
\(|P_m(C)-P(C)|=P_m(C)\leq1\). Hence \(\Delta_m=1\).

\emph{Item~3.} Let \(4(2m+1)e^{-m/32}<1\); such \(m\) exist because the
left side tends to zero. If \(D_m\) were completion-measurable, item~1 of
\cref{thm:p15-w02b-p15-065-outer-uniform} with \(\varepsilon=1/2\) and
\(\Pi_{\mathcal C}(2m)\leq2m+1\) would give
\(P^{*}(\Delta_m>1/2)<1\). That item uses only \(D_m\) for this \(m\), as
noted in its proof. But \(\{\Delta_m>1/2\}\) is the whole sample space,
of outer probability one. So \(D_m\) is not completion-measurable, and the
class is not symmetrization-\allowbreak measurable. By
\cref{thm:p15-w02b-p15-065-core-sufficient} it then has no countable core.
Directly, a countable family of countable sets has countable union, so a
point outside it cannot be approximated.

\emph{Item~4.} Let \(\widehat h_y(x)=\mathbf1\{x\prec y\}\) and
\(Y\equiv0\). Its risk is the probability of the countable set
\(\{x:x\prec y\}\), which is zero. For validation inputs
\(x_1,\ldots,x_v\), \(\widehat R_v(\widehat h_y)\) is the fraction of
indices with \(x_k\prec y\). For \(0\leq u<1\) the bad event \(F\) contains
\(G=\{(y,x_1,\ldots,x_v):x_k\prec y\text{ for all }k\}\). For fixed \(y\),
\(F\) is contained in \(\bigcup_k\{x_k\prec y\}\), a countable union of sets
whose \(k\)-th coordinate lies in a countable set. Its validation
probability is zero. Now let \(M\supseteq G\) be a measurable subset of
\([0,1]^{1+v}\). Its complement lies in
\(\{(y,x):y\preceq x_k\text{ for some }k\}\). For fixed \(x\), the
\(y\)-section of the complement lies in the countable set
\(\bigcup_kI_{x_k}\) and is null. By Tonelli's theorem the complement of
\(M\) is null, so \(M\) has probability one. Hence the outer probability of
\(G\), and of \(F\), is one. Each \(y\)-section of \(G\) lies in a null set,
so the iterated outer expectation of \(\mathbf1_G\) in
\cref{thm:p15-w02b-p15-065-outer-fubini}, inner over the validation
coordinates, is zero, while the joint outer expectation is one.
\end{proof}
\end{annexitem}

\begin{annexitem}{annex:p15-w02b-p15-065-real-uniform}{thm:p15-w02b-p15-065-real-uniform}{The layer-cake representation}
\begin{proof}
\emph{Item~1.} For \(g:Z\to[0,1]\) and each \(z\),
\(g(z)=\int_0^1\mathbf1\{g(z)>s\}\,ds\), since the integrand is the
indicator of \((0,g(z))\) within \((0,1)\). The set
\(\{(z,s):g(z)-s>0\}\) is product-measurable, so Tonelli's theorem gives
\(Pg=\int_0^1P\{g>s\}\,ds\); for \(P_m\) this is a finite sum. Hence
\[
 |P_mg-Pg|=\Bigl|\int_0^1\bigl(P_m\{g>s\}-P\{g>s\}\bigr)ds\Bigr|
 \leq\sup_{0<s<1}|P_m\{g>s\}-P\{g>s\}|,
\]
and each \(\{g>s\}\) is a member of \(\mathcal L_{\mathcal G}\).

\emph{Item~2.} By item~1 the bad event lies in
\(\{\Delta_m(\mathcal L_{\mathcal G})>\varepsilon\}\). Apply monotonicity of
\(P^{*}\) and item~1 of \cref{thm:p15-w02b-p15-065-outer-uniform}.

\emph{Item~3.} For \(\tilde g=(g-a)/(b-a)\),
\(|P_mg-Pg|=(b-a)|P_m\tilde g-P\tilde g|\) and
\(\{\tilde g>s\}=\{g>a+s(b-a)\}\). Apply item~2 with
\(\varepsilon/(b-a)\in(0,1)\).

\emph{Item~4.} If \(P_mg=0\), then \(g(Z_i)=0\) for all \(i\) because
\(g\geq0\), so \(P_m\{g>s\}=0\) for every \(s>0\). If moreover
\(Pg=\int_0^1P\{g>s\}ds>\varepsilon\), some \(s\in(0,1)\) has
\(P\{g>s\}>\varepsilon\). The event therefore lies in the event of item~2
of \cref{thm:p15-w02b-p15-065-outer-uniform} for \(\mathcal L_{\mathcal G}\).

\emph{Item~5.} As in item~3 of
\cref{annex:p15-w02b-p15-065-outer-learning}, the computation of
\cref{annex:p15.vcl.samples} bounds the right-hand sides of items~2 and~4 by
\(\delta\) at the displayed sizes, using only Sauer's bound with
\(d=d_{\mathcal G}\).
\end{proof}
\end{annexitem}

\begin{annexitem}{annex:p15-w02b-p15-065-real-learning}{thm:p15-w02b-p15-065-real-learning}{Infimum, approximate minimizers and a null exceptional set}
\begin{proof}
\emph{Item~1.} The empirical risks lie in \([0,1]\), so their infimum is a
real number. For \(\eta>0\) some member is within \(\eta\) of it, by the
definition of infimum. For the failure of exact minimization, take the
constant predictions \(\theta\in(0,1]\) and the loss \(\ell(\theta,y)=\theta\).
Every sample gives empirical risk \(\theta\), with infimum zero and no
minimizer.

\emph{Item~2.} On \(\{\sup_h|\widehat R_m(h)-R(h)|\leq\varepsilon\}\), for
every \(h\),
\(R(\widehat h)\leq\widehat R_m(\widehat h)+\varepsilon\leq\widehat R_m(h)
+\eta+\varepsilon\leq R(h)+2\varepsilon+\eta\); take the infimum over
\(h\). So the bad set lies in the event of item~2 of
\cref{thm:p15-w02b-p15-065-real-uniform}.

\emph{Item~3.} Since \(g_{h^{*}}\geq0\) and \(Pg_{h^{*}}=0\),
\(g_{h^{*}}=0\) almost everywhere. So
\(N=\bigcup_i\{g_{h^{*}}(Z_i)>0\}\) is a measurable null set, and off \(N\)
the member \(h^{*}\) has empirical risk zero. Off \(N\) the learner
therefore returns \(\widehat h\) with \(P_mg_{\widehat h}=0\), and
\(R(\widehat h)>\varepsilon\) places the sample in the event of item~4 of
\cref{thm:p15-w02b-p15-065-real-uniform}. Subadditivity gives the bound.
\end{proof}
\end{annexitem}

\begin{annexitem}{annex:p15-w02b-p15-065-noise}{thm:p15-w02b-p15-065-noise}{A zero-risk member bounds the floor}
\begin{proof}
\emph{Item~1.} Members of \(H\) are predictors, so
\(0\leq R^{*}\leq R(h)=0\). If \(R^{*}>0\), no member of any class has
risk zero, and the hypothesis of item~3 of
\cref{thm:p15-w02b-p15-065-real-learning} fails.

\emph{Item~2.} Let \(X\) be uniform on \([0,1]\), \(Y=X\), the loss
\(|y-a|\) on \([0,1]^2\), and \(H\) the constant predictions
\(\theta\in[0,1]\). The identity predictor has risk zero, so \(R^{*}=0\).
For a constant,
\(E|X-\theta|=\int_0^\theta(\theta-u)\,du+\int_\theta^1(u-\theta)\,du
=(\theta^2+(1-\theta)^2)/2\geq1/4\), so no member has risk zero.

\emph{Item~3.} Write
\(R(\widehat h)-R^{*}=(R(\widehat h)-R_H^{*})+(R_H^{*}-R^{*})\), as in
\cref{thm:p15-w02b-p15-064-decomposition} with the three-place loss
\(\ell(x,y,a)=\ell(a,y)\). On the event in question the first term is at
most \(2\varepsilon+\eta\) by
\cref{annex:p15-w02b-p15-065-real-learning}. The second term is defined
from \(H\), the loss and \(P\) alone.
\end{proof}
\end{annexitem}

\begin{annexitem}{annex:p15-w02b-p15-065-absolute-example}{thm:p15-w02b-p15-065-absolute-example}{Intervals, a rational core and the median}
\begin{proof}
\emph{Capacity.} The level sets are
\(L_{\theta,s}=\{(x,y):|y-\theta|>s\}\), \(\theta\in[0,1]\),
\(0<s<1\). Membership depends only on \(y\), so a shattered set has
distinct \(y\)-values. For two values \(a<b\) in \([0,1]\), with
\(c=(a+b)/2\) and \(r=(b-a)/2\in(0,1/2]\): \((\theta,s)=(c,r/2)\) contains
both; \((c,r)\) contains neither, since both lie at distance exactly \(r\);
\((b,r)\) contains only \(a\); \((a,r)\) contains only \(b\). For three
values \(a<b<c\), the complement \(\{y:|y-\theta|\leq s\}\) is an interval,
so it cannot contain \(a\) and \(c\) without \(b\). The pattern with only
\(b\) in the set is impossible. The dimension is two, and Sauer's bound
gives \(1+t+t(t-1)/2=(t^2+t+2)/2\).

\emph{Core.} Given \((\theta,s)\), choose rationals \(\theta_j\in[0,1]\)
with \(|\theta_j-\theta|<1/j\) and \(s_j\in(s+1/j,s+2/j)\), with \(j\) large
enough that \(s+2/j<1\). If \(|y-\theta|\leq s\), then
\(|y-\theta_j|<s+1/j<s_j\), so the approximants exclude \(y\). If
\(|y-\theta|>s\), then for large \(j\)
\(|y-\theta_j|\geq|y-\theta|-1/j>s+2/j>s_j\), so they include \(y\). The
rational pairs form a countable core, and
\cref{thm:p15-w02b-p15-065-core-sufficient} gives
symmetrization-\allowbreak measurability.

\emph{Risks.} The computation in \cref{annex:p15-w02b-p15-065-noise}
gives \(R(\theta)=(\theta^2+(1-\theta)^2)/2\), with derivative
\(2\theta-1\) and minimum \(1/4\) at \(\theta=1/2\). For any measurable
predictor \(f\) with values in \([0,1]\), independence and Tonelli's
theorem give \(R(f)=\int\phi(f(x))\,P_X(dx)\) with
\(\phi(a)=(a^2+(1-a)^2)/2\geq1/4\), so \(R^{*}=1/4=R_H^{*}\), and no
predictor has risk zero.

\emph{Sample size.} With \(d=2\), \(\varepsilon=0.1\), \(\delta=0.05\),
\(64/\varepsilon^2=6400\) and \(8/\delta=160\). Since
\(\ln6400=8.76405\ldots\) and \(\ln160=5.07517\ldots\), the bracket is
\(22.60328\ldots\) and the product is \(144660.99\ldots\), with ceiling
\(144661\). Item~5 of \cref{thm:p15-w02b-p15-065-real-uniform} and item~2
of \cref{thm:p15-w02b-p15-065-real-learning} with \(R_H^{*}=1/4\) give the
stated outer-probability bound.
\end{proof}
\end{annexitem}

\begin{annexitem}{annex:p15-w02b-p15-065-lower-realizable}{thm:p15-w02b-p15-065-lower-realizable}{Unseen points carry fair-coin labels}
\begin{proof}
Let the shattered points be \(x_0,\ldots,x_{d-1}\) and put
\(q=8\varepsilon\leq1\). Let \(\mu\) give mass \(1-q\) to \(x_0\) and
\(q/(d-1)\) to each \(x_i\), \(1\leq i\leq d-1\). For
\(b\in\{0,1\}^{d-1}\), let \(t_b(x_0)=0\), \(t_b(x_i)=b_i\), and let
\(P_b\) be the law of \((X,t_b(X))\) with \(X\sim\mu\). Shattering gives
\(h_b\in H\) equal to \(t_b\) on the \(d\) points, so \(R_{P_b}(h_b)=0\).
For any \(f\) put
\(R'_b(f)=\frac q{d-1}\sum_{i=1}^{d-1}\mathbf1\{f(x_i)\ne b_i\}\leq R_{P_b}(f)\).

Draw \(B\) uniform on \(\{0,1\}^{d-1}\) and, independently, indices
\(J_1,\ldots,J_m\) iid with \(\Pr(J=i)=\mu\{x_i\}\). The sample
\(S=((x_{J_k},t_B(x_{J_k})))_{k\leq m}\) has law \(P_B^{\otimes m}\) given
\(B\). Fix an index sequence \(j\) and let \(O\) be its set of indices. Given
\(J=j\), \(S\) is a function of \((B_k)_{k\in O}\). For \(i\notin O\),
\(A(S)(x_i)\) is therefore a function \(\psi\) of \((B_k)_{k\in O}\), which
is independent of the fair bit \(B_i\). So
\(\Pr(\psi\ne B_i)=1/2\). All the variables take finitely many values, so
these are finite sums. Averaging over \(j\),
\[
 E\,R'_B(A(S))\geq\frac q{d-1}\sum_{i=1}^{d-1}\tfrac12\Pr(i\notin O)
 =\frac q2\Bigl(1-\frac q{d-1}\Bigr)^m .
\]
Bernoulli's inequality \((1-x)^m\geq1-mx\) for \(x\in[0,1]\), proved by
induction on \(m\), and \(mq/(d-1)=8\varepsilon m/(d-1)\leq1/2\) give
\(E\,R'_B(A(S))\geq q/4=2\varepsilon\). For a randomized learner the same
holds conditionally on each internal draw \(u\). The integrand is a finite
sum of indicators measurable in \(u\), so integrating over \(u\) keeps the
bound. Since the average over \(b\) is at least \(2\varepsilon\), some \(b\)
has \(E_{P_b}R'_b(A(S))\geq2\varepsilon\). Because \(0\leq R'_b\leq q=8\varepsilon\),
\[
 2\varepsilon\leq\varepsilon\Pr(R'_b<\varepsilon)+8\varepsilon\Pr(R'_b\geq\varepsilon)
 \leq\varepsilon+8\varepsilon\Pr(R'_b\geq\varepsilon),
\]
so \(\Pr(R'_b\geq\varepsilon)\geq1/8\), and \(R_{P_b}\geq R'_b\) completes
the proof with \(P=P_b\), \(h^{*}=h_b\).
\end{proof}
\end{annexitem}

\begin{annexitem}{annex:p15-w02b-p15-065-lower-agnostic}{thm:p15-w02b-p15-065-lower-agnostic}{Biased coins and the affinity bound}
\begin{proof}
Let the shattered points be \(x_1,\ldots,x_d\) and \(\gamma=8\tau\leq1/2\).
For \(\sigma\in\{-1,1\}^d\), let \(P_\sigma\) make \(X\) uniform on the
\(d\) points and, given \(X=x_i\), make \(Y\) equal to one with probability
\(1/2+\sigma_i\gamma\). At \(x_i\) predicting \(1\) errs with probability
\(1/2-\sigma_i\gamma\), predicting \(0\) with \(1/2+\sigma_i\gamma\). The
rule \(f_\sigma(x_i)=\mathbf1\{\sigma_i=1\}\) attains the minimum
achievable risk \(1/2-\gamma\); shattering puts a member of \(H\) equal to
it on the support. So \(R_H^{*}=1/2-\gamma\), and for every \(f\)
\[
 R_{P_\sigma}(f)-R_H^{*}=\frac{2\gamma}d\sum_{i=1}^d\mathbf1\{f(x_i)\ne f_\sigma(x_i)\}\leq2\gamma .
\]

\emph{One coin.} For \(n\geq0\) let \(p_\pm\) be the laws on \(\{0,1\}^n\)
of \(n\) independent bits of mean \(1/2\pm\gamma\). With \(\sigma_i\)
uniform, any test \(\psi:\{0,1\}^n\to\{0,1\}\) errs with probability
\(\frac12\sum_L[p_+(L)\mathbf1\{\psi(L)=0\}+p_-(L)\mathbf1\{\psi(L)=1\}]
\geq\frac12\sum_L\min\{p_+,p_-\}\). By Cauchy--Schwarz,
\begin{align*}
 \Bigl(\sum_L\sqrt{p_+p_-}\Bigr)^2&=\Bigl(\sum_L\sqrt{\min\{p_+,p_-\}\max\{p_+,p_-\}}\Bigr)^2\\
 &\leq\sum_L\min\{p_+,p_-\}\sum_L\max\{p_+,p_-\}\leq2\sum_L\min\{p_+,p_-\}.
\end{align*}
The product form gives
\(\sum_L\sqrt{p_+p_-}=\bigl(2\sqrt{(1/2+\gamma)(1/2-\gamma)}\bigr)^n
=(1-4\gamma^2)^{n/2}\). Hence the error is at least
\(\frac14(1-4\gamma^2)^n\). For \(n=0\) this reads \(1/4\leq1/2\).

\emph{All coordinates.} Draw \(\sigma\) uniform, indices \(J_1,\ldots,J_m\)
iid uniform on \(\{1,\ldots,d\}\), and labels given \((\sigma,J)\)
independently with the stated means. The joint probability of
\((\sigma,J,Y)\) is \(2^{-d}d^{-m}\prod_k\beta(Y_k;\sigma_{J_k})\), where
\(\beta\) is the Bernoulli mass function. Given \(J\) it factors over the
points. So, given \(J\), the pair consisting of \(\sigma_i\) and the
\(n_i=\#\{k:J_k=i\}\) labels at \(x_i\) is independent of everything else,
\(W_i\). Now \(A(S)(x_i)\) is a function of those labels, \(W_i\) and
\(J\). For each fixed value of \((J,W_i)\) it is a test of \(\sigma_i\)
from \(n_i\) bits, with error at least \(\frac14(1-4\gamma^2)^{n_i}\), and
the error with respect to \(f_\sigma(x_i)\) is the error of that test.
The index count \(n_i\) is binomial with parameters \(m\) and \(1/d\), and its
generating function gives
\(E(1-4\gamma^2)^{n_i}=(1-4\gamma^2/d)^m\geq1-4\gamma^2m/d\), by
Bernoulli's inequality. Therefore
\[
 E\bigl[R_{P_\sigma}(A(S))-R_H^{*}\bigr]\geq\frac{2\gamma}d\cdot d\cdot
 \frac14\Bigl(1-\frac{4\gamma^2m}d\Bigr)\geq\frac\gamma4
\]
whenever \(m\leq d/(8\gamma^2)=d/(512\tau^2)\). A randomized learner is
handled by conditioning on its internal draw, as in
\cref{annex:p15-w02b-p15-065-lower-realizable}. Some \(\sigma\) attains at
least the average. For it, with \(V\) the excess risk, \(0\leq V\leq2\gamma\)
and \(EV\geq\gamma/4\) give
\(\gamma/4\leq\gamma/8+2\gamma\Pr(V\geq\gamma/8)\), so
\(\Pr(V\geq\gamma/8)\geq1/16\). Finally \(\gamma/8=\tau\).
\end{proof}
\end{annexitem}

\begin{annexitem}{annex:p15-w02b-p15-065-lower-infinite}{thm:p15-w02b-p15-065-lower-infinite}{Choosing the shattered set after the sample size}
\begin{proof}
Given \(m\) and \(\varepsilon\), choose an integer \(d\geq2\) with
\(d-1\geq16\varepsilon m\). Infinite VC dimension gives shattered finite
sets of every size, and every subset of a shattered set is shattered, so
some set of exactly \(d\) points is shattered. Then
\(m\leq(d-1)/(16\varepsilon)\), and
\cref{thm:p15-w02b-p15-065-lower-realizable} gives the realizable law.
A guarantee for every law would in particular be one for every realizable
law.
\end{proof}
\end{annexitem}

\begin{annexitem}{annex:p15-w02b-p15-065-noise-relations}{thm:p15-w02b-p15-065-noise-relations}{From the conditional label probability to the variance condition}
\begin{proof}
Throughout, \cref{thm:p15-w02b-p15-064-classification-bayes} with labels
\(\{0,1\}\) gives, for every measurable predictor \(f\),
\begin{equation*}
 R(f)-R^{*}=E\bigl[|2\eta(X)-1|\,\mathbf1\{f(X)\ne f_{\mathrm{Bayes}}(X)\}\bigr],
\end{equation*}
and characterizes the Bayes predictors as those agreeing with
\(f_{\mathrm{Bayes}}\) almost surely on \(\{\eta(X)\ne1/2\}\).

\emph{Item~1.} If \(Y=f^{\circ}(X)\) almost surely, then
\(\mathbf1\{Y=1\}=\mathbf1\{f^{\circ}(X)=1\}\) almost surely, and the right
side is \(\sigma(X)\)-measurable. Hence
\(\eta(X)=P(Y=1\mid X)=\mathbf1\{f^{\circ}(X)=1\}\in\{0,1\}\) almost
surely, and \(|2\eta(X)-1|=1\) almost surely.

\emph{Item~2.} Let \(0<t\leq1\). If \(t<\gamma\), the event
\(\{|2\eta(X)-1|\leq t\}\) is contained in \(\{|2\eta(X)-1|<\gamma\}\),
which is null, so its probability is \(0\leq\gamma^{-a}t^{a}\). If
\(t\geq\gamma\), its probability is at most \(1\leq(t/\gamma)^{a}\).

\emph{Item~3.} Under the bounded-noise condition,
\(|2\eta(X)-1|\geq\gamma>0\) almost surely, so \(\eta(X)\ne1/2\) almost
surely. Under the margin condition,
\(\{\eta(X)=1/2\}\subseteq\{|2\eta(X)-1|\leq t\}\) for every
\(0<t\leq1\), so \(P\{\eta(X)=1/2\}\leq Ct^{a}\) for all such \(t\), and
letting \(t\to0\) gives zero. In both cases \(\{\eta(X)\ne1/2\}\) has
probability one, and the characterization of Bayes predictors gives
\(f^{*}(X)=f_{\mathrm{Bayes}}(X)\) almost surely.

\emph{Item~4.} Let \(f^{*}\) be a Bayes predictor and \(f\) measurable.
By item~3, \(\mathbf1\{f(X)\ne f_{\mathrm{Bayes}}(X)\}=
\mathbf1\{f(X)\ne f^{*}(X)\}\) almost surely, and \(R(f^{*})=R^{*}\) is
the minimum over \(\mathcal F_{\mathrm{all}}\). The displayed identity and
\(|2\eta(X)-1|\geq\gamma\) almost surely give
\(R(f)-R(f^{*})\geq\gamma P\{f(X)\ne f^{*}(X)\}\), which is the variance
condition with \(\kappa=1\) and \(B=1/\gamma\).

\emph{Item~5.} Let \(f^{*}\) and \(f\) be as in item~4, and put
\(p=P\{f(X)\ne f^{*}(X)\}\) and \(\Delta=R(f)-R(f^{*})\). As in item~4,
\(\Delta=E[|2\eta(X)-1|\mathbf1\{f(X)\ne f^{*}(X)\}]\). If \(p=0\) the
inequality holds, since its right side is nonnegative. Let \(p>0\). For
\(0<t\leq1\), keeping only the part of the expectation where
\(|2\eta(X)-1|>t\),
\[
 \Delta\geq tP\{f(X)\ne f^{*}(X),\ |2\eta(X)-1|>t\}
 \geq t\bigl(p-P\{|2\eta(X)-1|\leq t\}\bigr)\geq t\bigl(p-Ct^{a}\bigr).
\]
Choose \(t=(p/(2C))^{1/a}\). Because \(0<p\leq1\leq C\), this \(t\) lies in
\((0,1)\), and \(Ct^{a}=p/2\). Hence
\[
 \Delta\geq\frac{tp}2=\frac{p^{(1+a)/a}}{2(2C)^{1/a}},
 \qquad\text{so}\qquad
 p\leq\bigl(2(2C)^{1/a}\Delta\bigr)^{a/(1+a)}
 =2^{a/(1+a)}(2C)^{1/(1+a)}\Delta^{a/(1+a)}.
\]
Since \(2^{a/(1+a)}2^{1/(1+a)}=2\), the constant is \(2C^{1/(1+a)}\), and
the exponent is \(\kappa=a/(1+a)\in(0,1)\).

\emph{Item~6.} Every probability is at most \(1=1\cdot\Delta^{0}\) under
the convention \(0^0=1\).

\emph{Item~7.} Let \(N=\{x:\eta(x)=1/2\}\), a measurable set with
\(P\{X\in N\}>0\), and let \(f^{*}\) be a Bayes predictor. Define \(f\) to
equal \(1-f^{*}\) on \(N\) and \(f^{*}\) off \(N\); it is measurable. It
agrees with \(f^{*}\), hence with \(f_{\mathrm{Bayes}}\), almost surely on
\(\{\eta(X)\ne1/2\}\), so it is a Bayes predictor and
\(R(f)-R(f^{*})=0\). For \(\kappa>0\) the right side of the variance
condition is \(B\cdot0^{\kappa}=0\), while
\(P\{f(X)\ne f^{*}(X)\}=P\{X\in N\}>0\).

\emph{Item~8.} If \(f^{*}\in\mathcal F\), then \(R(f^{*})=R^{*}\leq R(f)\)
for every \(f\in\mathcal F\), so \(f^{*}\) minimizes the risk over
\(\mathcal F\). The required inequality for each \(f\in\mathcal F\) is one
of the inequalities already assumed for \(\mathcal F_{\mathrm{all}}\).
\end{proof}
\end{annexitem}

\begin{annexitem}{annex:p15-w02b-p15-065-noise-rate}{thm:p15-w02b-p15-065-noise-rate}{Bernstein's form applied to loss differences}
\begin{proof}
\emph{Loss differences.} Fix \(h\in H\) and put
\[
 \xi_i=\mathbf1\{h(X_i)\ne Y_i\}-\mathbf1\{h^{\circ}(X_i)\ne Y_i\}.
\]
If \(h(X_i)=h^{\circ}(X_i)\) the two indicators coincide and
\(\xi_i=0\). If they differ, exactly one of the two binary predictions
equals \(Y_i\), so \(\xi_i\in\{-1,1\}\). Hence
\(\xi_i^2=\mathbf1\{h(X_i)\ne h^{\circ}(X_i)\}\), and with
\(v_h=P\{h(X)\ne h^{\circ}(X)\}\) and
\(\Delta_h=R(h)-R(h^{\circ})=R(h)-R_H^{*}\geq0\),
\[
 E\xi_i=\Delta_h,\qquad E\xi_i^2=v_h,\qquad
 m^{-1}\textstyle\sum_i\xi_i=\widehat R_m(h)-\widehat R_m(h^{\circ})
 =:\widehat\Delta_h .
\]
The variables \(U_i=\Delta_h-\xi_i\) are independent, have mean zero,
satisfy \(U_i\leq\Delta_h+1\leq2\), and have
\(EU_i^2=\operatorname{Var}\xi_i\leq v_h\).

\emph{One member.} Since \(r\geq1\) and \(0<\delta<1\), \(L>0\). Put
\[
 t_h=\sqrt{\frac{2v_hL}m}+\frac{4L}{3m}>0 .
\]
We show \(P\{\overline U\geq t_h\}\leq\delta/r=e^{-L}\), where
\(\overline U=m^{-1}\sum_iU_i=\Delta_h-\widehat\Delta_h\). If
\(\operatorname{Var}\xi_i=0\), then \(U_i=0\) almost surely and the
probability is zero. Otherwise \cref{thm:p15-variance-tail-chain}, in the
setting of \cref{def:p15-upper-variance-setting} with \(n=m\), \(b=2\) and
\(\sigma^2=\operatorname{Var}\xi_1>0\), gives in its Bernstein form, valid
for every \(t>0\),
\[
 P\{\overline U\geq t\}\leq\exp\Bigl\{-\frac{mt^2}{2(\sigma^2+2t/3)}\Bigr\}
 \leq\exp\Bigl\{-\frac{mt^2}{2(v_h+2t/3)}\Bigr\},
\]
the second step because the exponent decreases in absolute value as the
variance term grows. The last bound is at most \(e^{-L}\) exactly when
\(q(t)=t^2-\frac{4L}{3m}t-\frac{2v_hL}m\geq0\). The larger root of \(q\) is
\[
 t_+=\frac{2L}{3m}+\sqrt{\frac{4L^2}{9m^2}+\frac{2v_hL}m}
 \leq\frac{2L}{3m}+\frac{2L}{3m}+\sqrt{\frac{2v_hL}m}=t_h,
\]
using \(\sqrt{x+y}\leq\sqrt x+\sqrt y\), and \(q\geq0\) on
\([t_+,\infty)\). So \(q(t_h)\geq0\) and the claim follows.

\emph{All members.} Each event \(\{\overline U\geq t_h\}\) is measurable,
because \(h\) and \(h^{\circ}\) are measurable. Let \(E\) be the
intersection over the \(r\) members of the complements. By the union
bound, \(P(E)\geq1-r\cdot\delta/r=1-\delta\). On \(E\), for every
\(h\in H\), the variance condition \(v_h\leq B\Delta_h^{\kappa}\) gives
\[
 \Delta_h<\widehat\Delta_h+\sqrt{\frac{2v_hL}m}+\frac{4L}{3m}
 \leq\widehat\Delta_h+\sqrt{\frac{2B\Delta_h^{\kappa}L}m}+\frac{4L}{3m},
\]
which is the displayed inequality. The event \(E\) is defined before any
selection, so the conclusions below hold for every member, and hence for
any member a selection rule returns, measurable or not.

\emph{Item~1.} Let \(h\) satisfy \(\widehat\Delta_h\leq0\), write
\(\Delta=\Delta_h\) and \(S=\sqrt{2B\Delta^{\kappa}L/m}\). On \(E\),
\(\Delta\leq S+4L/(3m)\). If \(\Delta\leq8L/(3m)\) there is nothing to
prove. Otherwise \(4L/(3m)<\Delta/2\), so \(\Delta<2S\) and \(\Delta>0\);
squaring, \(\Delta^2<8B\Delta^{\kappa}L/m\), so
\(\Delta^{2-\kappa}<8BL/m\), and since \(2-\kappa\geq1\),
\(\Delta<(8BL/m)^{1/(2-\kappa)}\). An exact empirical minimizer has
\(\widehat R_m(h)\leq\widehat R_m(h^{\circ})\).

\emph{Item~2.} An \(\vartheta\)-approximate minimizer has
\(\widehat R_m(h)\leq\min_{g\in H}\widehat R_m(g)+\vartheta
\leq\widehat R_m(h^{\circ})+\vartheta\), so \(\widehat\Delta_h\leq\vartheta\) and
\(\Delta\leq\vartheta+S+4L/(3m)\leq3\max\{\vartheta,S,4L/(3m)\}\). If the maximum
is \(\vartheta\), then \(\Delta\leq3\vartheta\); if it is \(4L/(3m)\), then
\(\Delta\leq4L/m\). If it is \(S\) and \(\Delta>0\), then
\(\Delta^2\leq9S^2=18B\Delta^{\kappa}L/m\), so
\(\Delta\leq(18BL/m)^{1/(2-\kappa)}\); if \(\Delta=0\) every bound holds.

\emph{Bayes reference.} By the definition of the variance condition,
\(h^{\circ}\) minimizes the risk over \(H\), so \(R_H^{*}=R(h^{\circ})\).
If \(h^{\circ}\) is a Bayes predictor this equals \(R^{*}\).
\end{proof}
\end{annexitem}

\begin{annexitem}{annex:p15-w02b-p15-065-noise-sizes}{thm:p15-w02b-p15-065-noise-sizes}{Solving the rate for the sample size}
\begin{proof}
If \(m\geq8BL/\varepsilon^{2-\kappa}\), then
\(8BL/m\leq\varepsilon^{2-\kappa}\), so
\((8BL/m)^{1/(2-\kappa)}\leq\varepsilon\). If
\(m\geq8L/(3\varepsilon)\), then \(8L/(3m)\leq\varepsilon\). Item~1 of
\cref{thm:p15-w02b-p15-065-noise-rate} then gives
\(R(\widehat h)-R_H^{*}\leq\varepsilon\) on an event of probability at
least \(1-\delta\), for every exact empirical minimizer \(\widehat h\).

Now let \(f^{*}\in H\) be a Bayes predictor and take \(h^{\circ}=f^{*}\);
it minimizes the risk over \(H\), and \(R_H^{*}=R^{*}\).

\emph{Item~1.} If \(Y=h(X)\) almost surely for some \(h\in H\), then
\(R(h)=0\), so \(R^{*}=0\) and \(h\) is a Bayes predictor; take
\(h^{\circ}=h\). Items~1, 4 and~8 of
\cref{thm:p15-w02b-p15-065-noise-relations} give the variance condition on
\(H\) with \(\kappa=1\) and \(B=1\), and the requirement becomes
\(m\geq\max\{8L/\varepsilon,8L/(3\varepsilon)\}=8L/\varepsilon\).

\emph{Item~2.} Items~4 and~8 give \(\kappa=1\) and \(B=1/\gamma\geq1\), so
\(8BL/\varepsilon=8L/(\gamma\varepsilon)\geq8L/(3\varepsilon)\).

\emph{Item~3.} Items~5 and~8 give \(\kappa=a/(1+a)\) and
\(B=2C^{1/(1+a)}\), so \(2-\kappa=(2+a)/(1+a)\) and
\(8BL/\varepsilon^{2-\kappa}=16C^{1/(1+a)}L\varepsilon^{-(2+a)/(1+a)}\).
Since \(C\geq1\), \(0<\varepsilon\leq1\) and \((2+a)/(1+a)\geq1\), this is
at least \(16L/\varepsilon\geq8L/(3\varepsilon)\).

\emph{Item~4.} Item~6 gives \(\kappa=0\) and \(B=1\) relative to
\(h^{\circ}\), and \(8L/\varepsilon^2\geq8L/(3\varepsilon)\) for
\(\varepsilon\leq1\).

Finally \((2+a)/(1+a)=1+1/(1+a)\), which decreases strictly in \(a\),
tends to \(2\) as \(a\to0\) and to \(1\) as \(a\to\infty\).
\end{proof}
\end{annexitem}

\begin{annexitem}{annex:p15-w02b-p15-065-noise-lower}{thm:p15-w02b-p15-065-noise-lower}{The agnostic construction has bounded noise}
\begin{proof}
Put \(\tau=\gamma/16\), so \(0<\tau\leq1/16\), and
\(d/(512\tau^2)=d/(2\gamma^2)\geq m\).
\Cref{thm:p15-w02b-p15-065-lower-agnostic} applies. Its proof
(\cref{annex:p15-w02b-p15-065-lower-agnostic}) produces, for the given
learner, one of the laws \(P_\sigma\), \(\sigma\in\{-1,1\}^d\), with
\(P_\sigma^{\otimes m}\{R_{P_\sigma}(A(S))-R_H^{*}\geq\tau\}\geq1/16\).
Under \(P_\sigma\), \(X\) is uniform on the shattered points
\(x_1,\ldots,x_d\), and \(Y\) given \(X=x_i\) equals one with probability
\(1/2+8\tau\sigma_i\). Singletons are measurable, so
\[
 \eta=\sum_{i=1}^d\Bigl(\frac12+8\tau\sigma_i\Bigr)\mathbf1_{\{x_i\}}
 +\frac12\,\mathbf1_{\mathcal X\setminus\{x_1,\ldots,x_d\}}
\]
is measurable, and \(\eta(X)\) is a version of \(P(Y=1\mid X)\) because
each support point has positive probability. Almost surely
\(|2\eta(X)-1|=16\tau=\gamma\), so the bounded-noise condition holds with
margin \(\gamma\). The same proof shows that the rule \(f_\sigma\) with
\(f_\sigma(x_i)=\mathbf1\{\sigma_i=1\}\) attains the minimum achievable
risk, and that some \(h\in H\) agrees with it on the support; then
\(R_{P_\sigma}(h)=R_{P_\sigma}(f_\sigma)=R^{*}\), so \(h\) is a Bayes
predictor and \(R_H^{*}=R^{*}\). Substituting \(\tau=\gamma/16\) gives the
display, and \(d/(2\gamma^2)=d/(32\gamma\cdot\gamma/16)\) gives the last
sentence.
\end{proof}
\end{annexitem}

\begin{annexitem}{annex:p15-w02b-p15-065-noise-example}{thm:p15-w02b-p15-065-noise-example}{Interval lengths and three sample sizes}
\begin{proof}
\emph{Bayes predictor and noise.} \(f_{\mathrm{Bayes}}(x)=
\mathbf1\{\eta(x)\geq1/2\}=\mathbf1\{x\geq3/10\}=h_{30}(x)\). Since
\(\min\{\eta,1-\eta\}=1/10\) everywhere,
\cref{thm:p15-w02b-p15-064-classification-bayes} gives \(R^{*}=1/10\).
Also \(|2\eta(x)-1|=4/5\) for every \(x\), so the bounded-noise condition
holds with \(\gamma=4/5\).

\emph{Variance condition with equality.} For \(j\in\{0,\ldots,100\}\),
\(h_j\) and \(h_{30}\) differ exactly on the interval between \(j/100\)
and \(3/10\), closed at the smaller end and open at the larger, of length
\(|j-30|/100\). The identity in
\cref{thm:p15-w02b-p15-064-classification-bayes} gives
\(R(h_j)-R^{*}=\tfrac45\cdot|j-30|/100\). Hence
\(P\{h_j(X)\ne h_{30}(X)\}=\tfrac54(R(h_j)-R(h_{30}))\) for every \(j\),
and \(h_{30}\) minimizes the risk over \(H\).

\emph{Sizes.} With \(r=101\), \(\delta=1/20\) and \(\varepsilon=1/100\),
\(L=\ln(101\cdot20)=\ln2020=7.610852\ldots\). Item~2 of
\cref{thm:p15-w02b-p15-065-noise-sizes} with \(\gamma=4/5\) requires
\(m\geq8L/(\gamma\varepsilon)=1000L=7610.852\ldots\), so the smallest
admissible integer is \(7611\). The size condition of
\cref{thm:p15-conc-finite-learning} is
\(n\geq\lceil2\ln(2r/\delta)/\varepsilon^2\rceil
=\lceil20000\ln4040\rceil\), and \(\ln4040=8.30399997\ldots\), so
\(20000\ln4040=166079.9994\ldots\) and the requirement is
\(166080\). Item~4 of \cref{thm:p15-w02b-p15-065-noise-sizes} requires
\(m\geq80000L=608868.22\ldots\), that is \(m\geq608869\).

\emph{Conclusion.} For \(m=7611\) the corollary gives
\(R(\widehat h)-R^{*}\leq1/100\), that is \(R(\widehat h)\leq0.11\), with
probability at least \(1-1/20=0.95\), for every exact empirical
minimizer.
\end{proof}
\end{annexitem}

\begin{annexitem}{annex:p15-w02b-p15-065-ternary-core}{thm:p15-w02b-p15-065-ternary-core}{Patterns on the sample}
\begin{proof}
\emph{Item~1.} Fix \((x,y)\). If \(h(x)=h^{\circ}(x)\), the two indicators
in \(\xi_h\) coincide, so \(\xi_h(x,y)=0=\chi_h(x,y)\). If
\(h(x)\ne h^{\circ}(x)\), exactly one of the two binary values equals
\(y\), so \(\xi_h(x,y)\in\{-1,1\}\) and
\(\xi_h(x,y)^2=|\xi_h(x,y)|=1=\chi_h(x,y)\). Integrating,
\(P\chi_h=P\{h(X)\ne h^{\circ}(X)\}\), and by linearity
\(P\xi_h=P\{h(X)\ne Y\}-P\{h^{\circ}(X)\ne Y\}=R(h)-R(h^{\circ})\).

\emph{Item~2.} Fix an ordered tuple \(((x_1,y_1),\ldots,(x_t,y_t))\). For
each \(k\), \(\xi_h(x_k,y_k)\) and \(-\chi_h(x_k,y_k)\) are determined by
\(h(x_k)\), because \(y_k\) and \(h^{\circ}(x_k)\) do not depend on \(h\).
So the vector of values of \(\xi_h\), and the vector of values of
\(-\chi_h\), on the tuple are functions of the vector
\[
 (h(x_1),\ldots,h(x_t))=(\mathbf1_{\{h=1\}}(x_1),\ldots,\mathbf1_{\{h=1\}}(x_t)).
\]
By the definition of \(\Pi_H\) the latter takes at most \(\Pi_H(t)\)
values as \(h\) ranges over \(H\). Hence each of the two families leaves
at most \(\Pi_H(t)\) vectors, and \(\mathcal T_H\) at most \(2\Pi_H(t)\).

\emph{Item~3.} Let \(H_0\subseteq H\) be countable such that each
\(\mathbf1_{\{h=1\}}\) is a pointwise limit of indicators
\(\mathbf1_{\{h_j=1\}}\) with \(h_j\in H_0\); since the predictors are
binary, this says \(h_j\to h\) pointwise. Fix \(w\in Z^{2m}\) and
\(h\in H\). A convergent binary sequence is eventually constant, so for
all large \(j\), \(h_j\) agrees with \(h\) at each of the \(2m\) inputs
of \(w\). By item~2, \(\xi_{h_j}\) and \(\xi_h\) then take the same values
at every coordinate of \(w\), and so do \(-\chi_{h_j}\) and \(-\chi_h\).
For \(g\in\mathcal T\) put
\[
 \varphi_g(w)=\frac{m(P'_mg-P_mg)(w)}{\sqrt{n_g(w)}}\ \text{if }n_g(w)>0,
 \qquad\varphi_g(w)=0\ \text{if }n_g(w)=0 .
\]
Then
\[
 J_m(w)=\max\Bigl\{0,\ \sup_{g\in\mathcal T_H}\varphi_g(w)\Bigr\},
\]
and \(\varphi_g(w)\) depends on \(g\) only through its values on the
coordinates of \(w\). Each \(\varphi_g(w)\) with \(g\in\mathcal T_H\)
therefore equals some \(\varphi_{g'}(w)\) with \(g'\) in the countable
class
\[
 \mathcal T_{H_0}=\{\xi_{h'}:h'\in H_0\}\cup\{-\chi_{h'}:h'\in H_0\}\subseteq\mathcal T_H,
\]
so the supremum over \(\mathcal T_H\) equals that over \(\mathcal T_{H_0}\).
For measurable \(g\), the maps \(w\mapsto n_g(w)\) and
\(w\mapsto m(P'_mg-P_mg)(w)\) are finite sums of measurable functions of
single coordinates, so \(\varphi_g\) is \(\mathcal A^{\otimes2m}\)-measurable.
A countable supremum of measurable functions is measurable, and so is its
maximum with \(0\). A measurable function is measurable for the
completion under any law.
\end{proof}
\end{annexitem}

\begin{annexitem}{annex:p15-w02b-p15-065-relative}{thm:p15-w02b-p15-065-relative}{Normalized swaps}
\begin{proof}
If \(2\Pi_{\mathcal T}(2m)e^{-L}\geq1\) there is nothing to prove. Assume
otherwise. Since \(\Pi_{\mathcal T}(2m)\geq1\), then \(e^{-L}<1/2\), so
\(L>\ln2>1/5\). Let \(A\subseteq Z^m\) be the set in the statement.

\emph{The ghost sample.} For each \(z\in A\) choose \(g=g_z\in\mathcal T\)
with \(Pg-P_mg(z)>t\), where
\[
 v=P|g|,\qquad \widehat p=P_m|g|(z),\qquad M=\max\{v,\widehat p\},\qquad
 t=\sqrt{80ML/m};
\]
the choice is made separately for each \(z\), with no measurability of
\(z\mapsto g_z\). If \(M=0\), then \(v=0\), so \(g=0\) almost surely and
\(Pg=0\), and \(\widehat p=0\), so \(P_mg(z)=0\); this contradicts
\(Pg-P_mg(z)>0\). Hence \(M>0\) and \(t>0\). Let \(Z'_1,\ldots,Z'_m\) be
iid with law \(P\), and let
\[
 G_z=\bigl\{z':\ P'_mg\geq Pg-t/2,\ P'_m|g|\leq4v\bigr\},
\]
a measurable set because \(g\) is measurable. Since \(g^2=|g|\),
\(P'_mg\) has mean \(Pg\) and variance at most \(Pg^2/m=v/m\).
Chebyshev's inequality gives
\[
 P^{\otimes m}\{P'_mg<Pg-t/2\}\leq\frac{4v}{mt^2}=\frac{4v}{80ML}
 \leq\frac1{20L}<\frac14 .
\]
If \(v>0\), Markov's inequality gives
\(P^{\otimes m}\{P'_m|g|>4v\}\leq v/(4v)=1/4\); if \(v=0\), then
\(P'_m|g|=0\) almost surely. Hence \(P^{\otimes m}(G_z)\geq1/2\).

Let \(z'\in G_z\) and \(w=(z,z')\). Then
\[
 m(P'_mg-P_mg)(w)\geq m\bigl(Pg-t/2-P_mg(z)\bigr)>mt/2>0,
\]
so \(g\) is not zero at some coordinate of \(w\) and \(n_g(w)>0\).
Moreover \(n_g(w)=m(\widehat p+P'_m|g|)\leq m(\widehat p+4v)\leq5mM\).
Therefore
\[
 \frac{m(P'_mg-P_mg)(w)}{\sqrt{n_g(w)}}>\frac{mt}{2\sqrt{5mM}}
 =\frac12\sqrt{\frac{80mML}{5mM}}=2\sqrt L,
\]
and \(J_m(w)>2\sqrt L\).

\emph{From the sample to the pooled sample.} Let
\(f=\mathbf1\{J_m>2\sqrt L\}\) on \(Z^{2m}\). For \(z\in A\),
\(f(z,\cdot)\geq\mathbf1_{G_z}\), so \(E^{*}f(z,\cdot)\geq1/2\) by the
elementary properties in \cref{annex:p15-w02b-p15-065-outer-fubini}. Thus
\(\mathbf1_A(z)\leq2E^{*}f(z,\cdot)\) for every \(z\in Z^m\), and
monotonicity, homogeneity and the product inequality of
\cref{thm:p15-w02b-p15-065-outer-fubini} give
\[
 P^{*}(A)\leq2E^{*}_{P^{\otimes m}}\bigl[E^{*}_{P^{\otimes m}}f(z,\cdot)\bigr]
 \leq2E^{*}_{P^{\otimes2m}}f=2\,\overline{P^{\otimes2m}}(W),
\]
where \(W=\{J_m>2\sqrt L\}\) is completion-measurable by hypothesis.

\emph{Swaps.} The swap maps \(T_\sigma\) preserve the completed law and
completion-measurability, as shown at the start of
\cref{annex:p15-w02b-p15-065-outer-uniform}. Hence
\[
 \overline{P^{\otimes2m}}(W)=2^{-m}\sum_\sigma\overline{P^{\otimes2m}}(T_\sigma^{-1}W)
 =\int\Pr_\sigma\bigl(J_m(T_\sigma w)>2\sqrt L\bigr)\,\overline{P^{\otimes2m}}(dw),
\]
with \(\sigma\) uniform on \(\{-1,1\}^m\). Fix \(w\). For \(g\in\mathcal T\)
put \(c_i=g(z'_i)-g(z_i)\). The first half of \(T_\sigma w\) holds \(z_i\)
where \(\sigma_i=1\) and \(z'_i\) where \(\sigma_i=-1\), so
\[
 m(P'_mg-P_mg)(T_\sigma w)=\sum_i\sigma_ic_i,\qquad n_g(T_\sigma w)=n_g(w).
\]
Hence \(J_m(T_\sigma w)>2\sqrt L\) iff some \(g\in\mathcal T\) with
\(n_g(w)>0\) has \(\sum_i\sigma_ic_i>2\sqrt{Ln_g(w)}\), and this event for
\(g\) depends on \(g\) only through the vector of its values on the
\(2m\) coordinates of \(w\). At most \(\Pi_{\mathcal T}(2m)\) such vectors
occur. Fix one, with \(n=n_g(w)>0\) and \(s=\sum_ic_i^2\). If \(s=0\) the
sum is zero and the event is empty. Otherwise, as in
\cref{annex:p15-w02b-p15-065-outer-uniform}, \(\cosh u\leq e^{u^2/2}\) for
real \(u\), independence of the signs and Markov's inequality give, for
\(\lambda>0\) and \(u_0=2\sqrt{Ln}\),
\[
 \Pr_\sigma\Bigl(\sum_i\sigma_ic_i>u_0\Bigr)\leq e^{-\lambda u_0}\prod_i\cosh(\lambda c_i)
 \leq e^{-\lambda u_0+\lambda^2s/2},
\]
which at \(\lambda=u_0/s\) equals \(e^{-u_0^2/(2s)}\). Since
\((a-b)^2\leq2a^2+2b^2\) and \(g^2=|g|\),
\(s\leq2\sum_i(|g(z_i)|+|g(z'_i)|)=2n\). So the probability is at most
\(e^{-4Ln/(4n)}=e^{-L}\). A union over the vectors gives
\(\Pr_\sigma(J_m(T_\sigma w)>2\sqrt L)\leq\Pi_{\mathcal T}(2m)e^{-L}\) for
every \(w\), with no measurability needed since \(\sigma\) ranges over a
finite set. Integrating, \(\overline{P^{\otimes2m}}(W)\leq\Pi_{\mathcal T}(2m)e^{-L}\),
and \(P^{*}(A)\leq2\Pi_{\mathcal T}(2m)e^{-L}\).
\end{proof}
\end{annexitem}

\begin{annexitem}{annex:p15-w02b-p15-065-vc-noise-rate}{thm:p15-w02b-p15-065-vc-noise-rate}{Two applications}
\begin{proof}
\emph{The exceptional set.} Apply
\cref{thm:p15-w02b-p15-065-relative} to \(\mathcal T_H\) with the given
\(L\); its measurability hypothesis is assumed. By item~2 of
\cref{thm:p15-w02b-p15-065-ternary-core},
\[
 2\Pi_{\mathcal T_H}(2m)e^{-L}\leq4\Pi_H(2m)e^{-L}=\delta .
\]
So the set \(A\) of samples on which some \(g\in\mathcal T_H\) has
\[
 Pg-P_mg>\sqrt{80\max\{P|g|,P_m|g|\}L/m}
\]
has \(P^{*}(A)\leq\delta\).
Fix a sample outside \(A\) and \(h\in H\), and write
\[
 \Delta_h=R(h)-R_H^{*},\quad \widehat\Delta_h=\widehat R_m(h)-\widehat R_m(h^{\circ}),\quad
 v_h=P\chi_h,\quad p_h=P_m\chi_h .
\]
By item~1 of \cref{thm:p15-w02b-p15-065-ternary-core} and
\(R(h^{\circ})=R_H^{*}\), \(\Delta_h=P\xi_h\geq0\); also
\(\widehat\Delta_h=P_m\xi_h\), and \(|\xi_h|=|-\chi_h|=\chi_h\).

\emph{The disagreement frequency.} With \(g=-\chi_h\),
\(Pg-P_mg=p_h-v_h\leq\sqrt{80\max\{v_h,p_h\}L/m}\). If \(p_h\leq v_h\),
then \(\max\{v_h,p_h\}=v_h\). If \(p_h>v_h\), then
\(p_h-v_h\leq\sqrt{p_h\cdot80L/m}\leq p_h/2+40L/m\), by
\(\sqrt{xy}\leq(x+y)/2\), so \(p_h\leq2v_h+80L/m\). In both cases
\(\max\{v_h,p_h\}\leq2v_h+80L/m\).

\emph{The loss difference.} With \(g=\xi_h\), and
\(\sqrt{x+y}\leq\sqrt x+\sqrt y\),
\[
 \Delta_h-\widehat\Delta_h\leq\sqrt{\frac{80(2v_h+80L/m)L}m}
 \leq\sqrt{\frac{160v_hL}m}+\frac{80L}m .
\]
The variance condition \(v_h\leq B\Delta_h^{\kappa}\) gives the displayed
inequality, for every \(h\in H\) and every sample outside \(A\).

\emph{Item~1.} Let \(\widehat\Delta_h\leq0\), write \(\Delta=\Delta_h\) and
\(S=\sqrt{160B\Delta^{\kappa}L/m}\), so \(\Delta\leq S+80L/m\). If
\(\Delta\leq160L/m\) the bound holds. Otherwise \(80L/m<\Delta/2\), so
\(\Delta<2S\) and \(\Delta>0\); squaring,
\(\Delta^2<640B\Delta^{\kappa}L/m\), so \(\Delta^{2-\kappa}<640BL/m\), and
since \(2-\kappa\geq1\), \(\Delta<(640BL/m)^{1/(2-\kappa)}\). The
empirical risk takes values in \(\{0,1/m,\ldots,1\}\), so an exact
empirical minimizer exists, and it has
\(\widehat R_m(h)\leq\widehat R_m(h^{\circ})\).

\emph{Item~2.} A \(\vartheta\)-approximate minimizer has
\(\widehat R_m(h)\leq\inf_{g\in H}\widehat R_m(g)+\vartheta\leq
\widehat R_m(h^{\circ})+\vartheta\), so \(\widehat\Delta_h\leq\vartheta\) and
\(\Delta\leq\vartheta+S+80L/m\leq3\max\{\vartheta,S,80L/m\}\). If the
maximum is \(\vartheta\), \(\Delta\leq3\vartheta\); if it is \(80L/m\),
\(\Delta\leq240L/m\). If it is \(S\) and \(\Delta>0\), then
\(\Delta^2\leq9S^2=1440B\Delta^{\kappa}L/m\), so
\(\Delta\leq(1440BL/m)^{1/(2-\kappa)}\); if \(\Delta=0\) every bound holds.

\emph{Growth.} If \(1\leq d\leq2m\), \cref{thm:p15-vcl-sauer} with
\(t=2m\) gives \(\Pi_H(2m)\leq(2em/d)^d\), so
\(L\leq d\ln(2em/d)+\ln(4/\delta)\).

\emph{Selection.} The set \(A\) is fixed before any member is selected,
and outside it the conclusions hold for every member of \(H\). They
therefore hold for whatever member a selection rule returns, measurable or
not.
\end{proof}
\end{annexitem}

\begin{annexitem}{annex:p15-w02b-p15-065-vc-noise-orders}{thm:p15-w02b-p15-065-vc-noise-orders}{Substitution}
\begin{proof}
A Bayes predictor \(f^{*}\in H\) has \(R(f^{*})=R^{*}\leq R(h)\) for every
measurable \(h\), so it minimizes the risk over \(H\) and
\(R_H^{*}=R^{*}\). Each item applies item~1 of
\cref{thm:p15-w02b-p15-065-vc-noise-rate} with \(h^{\circ}=f^{*}\) and the
exponent and constant supplied by
\cref{thm:p15-w02b-p15-065-noise-relations}, restricted to \(H\) by its
item~8.

\emph{Item~1.} If \(Y=h(X)\) almost surely for some \(h\in H\), then
\(R(h)=0\), so \(R^{*}=0\). Items~1 and~4 of the lemma give the variance
condition with \(\kappa=1\) and \(B=1\) relative to every Bayes predictor,
in particular \(f^{*}\). The bound is
\(\max\{640L/m,160L/m\}=640L/m\), and \(R(\widehat h)-R^{*}=R(\widehat h)\).

\emph{Item~2.} Item~4 gives \(\kappa=1\) and \(B=1/\gamma\geq1\), so the
bound is \(\max\{640L/(\gamma m),160L/m\}=640L/(\gamma m)\).

\emph{Item~3.} Item~5 gives \(\kappa=a/(1+a)\) and
\(B=2C^{1/(1+a)}\). Then \(2-\kappa=(2+a)/(1+a)\), so
\(1/(2-\kappa)=(1+a)/(2+a)\), and \(640B=1280C^{1/(1+a)}\).

\emph{Item~4.} Item~6 gives \(\kappa=0\) and \(B=1\) relative to the
minimizing member \(f^{*}\), so \(1/(2-\kappa)=1/2\).

\emph{Orders.} Under \(1\leq d\leq2m\), the bound on \(L\) is the
growth step of \cref{annex:p15-w02b-p15-065-vc-noise-rate}, and
\(d\ln(2em/d)=d\ln m+d\ln(2e/d)\).
\end{proof}
\end{annexitem}

\begin{annexitem}{annex:p15-w02b-p15-065-vc-noise-example}{thm:p15-w02b-p15-065-vc-noise-example}{Rational thresholds}
\begin{proof}
\emph{Core.} Let \(H_0=\{h_q:q\in[0,1]\cap\mathbb Q\}\). For \(s=0\),
\(h_0\in H_0\). For \(0<s\leq1\), choose rationals \(q_j\in[0,s)\) with
\(q_j\to s\). If \(x\geq s\), then \(x>q_j\) and \(h_{q_j}(x)=1=h_s(x)\)
for every \(j\). If \(x<s\), then \(q_j>x\) for all large \(j\), so
\(h_{q_j}(x)=0=h_s(x)\) eventually. So the indicators of
\(\{h_{q_j}=1\}\) converge pointwise to that of \(\{h_s=1\}\), and
item~3 of \cref{thm:p15-w02b-p15-065-ternary-core} gives the measurability
hypothesis of \cref{thm:p15-w02b-p15-065-vc-noise-rate} for every law and
every \(m\).

\emph{Growth.} The positive prediction-events are the intervals
\([s,1]\). Two points \(x<x'\) are not shattered: an interval \([s,1]\)
containing \(x\) contains \(x'\). So the VC dimension is at most one, and
\cref{thm:p15-vcl-sauer} gives \(\Pi_H(t)\leq\binom t0+\binom t1=t+1\).
Conversely, on \(t\) distinct points \(x_1<\cdots<x_t\) in \([0,1)\),
the thresholds \(s=x_1,\ldots,x_t\) and \(s=1\) give vectors with
exactly \(t,t-1,\ldots,1,0\) entries equal to one, so
\(\Pi_H(t)=t+1\).

\emph{Bayes predictor and variance condition.} As in
\cref{annex:p15-w02b-p15-065-noise-example}, which treats the finite
example, \(f_{\mathrm{Bayes}}=h_{3/10}\), \(R^{*}=1/10\), and the bounded-noise
condition holds with \(\gamma=4/5\). For \(s\in[0,1]\), \(h_s\) and
\(h_{3/10}\) differ exactly on the interval between \(s\) and \(3/10\),
closed at the smaller end and open at the larger, of length
\(|s-3/10|\). The identity in
\cref{thm:p15-w02b-p15-064-classification-bayes}, with
\(|2\eta-1|=4/5\) everywhere, gives \(R(h_s)-R^{*}=\tfrac45|s-3/10|\).
Hence \(P\{h_s(X)\ne h_{3/10}(X)\}=\tfrac54(R(h_s)-R(h_{3/10}))\) for every
\(s\), \(h_{3/10}\) minimizes the risk over \(H\), and the variance
condition holds with \(\kappa=1\) and \(B=5/4\).

\emph{Sample size.} With \(m=1550000\) and \(\delta=1/20\),
\(L=\ln(4\cdot3100001\cdot20)=\ln248000080=19.3289\ldots\). Numerically
\(e^{19.3}=2.409\ldots\times10^{8}\) and
\(e^{19.34}=2.507\ldots\times10^{8}\), so \(19.3<L<19.34\). Item~2 of
\cref{thm:p15-w02b-p15-065-vc-noise-orders} with \(\gamma=4/5\) gives
\[
 R(\widehat h)-R^{*}\leq\frac{640L}{\gamma m}=\frac{800L}{1550000}
 <\frac{800\cdot19.34}{1550000}=0.00998\ldots<0.01
\]
outside a set of outer probability at most \(1/20\), for every exact
empirical minimizer; hence \(R(\widehat h)<0.11\) there. The bound of
item~4 is at least
\[
 \sqrt{\frac{640L}{m}}>\sqrt{\frac{640\cdot19.3}{1550000}}
 =\sqrt{0.0079690\ldots}=0.08926\ldots>0.089 .
\]
\end{proof}
\end{annexitem}
